% SPDX-License-Identifier: Apache-2.0
% covers: ScanCtl, ScanTap
\datasheet{ScanCtl}{The scanout's registers: the frame's base, the bits
that show the scanout and ask two rows ahead, and the underflow and
stuck bits, read and cleared, with the address of a line that did not
come, the longest line and the count of late lines.}

\dsfacts{\code{lib/parts/src/scanout.rs}}%
{\code{txhdl\_parts::scanout::ScanCtl}}%
{\code{//docs:examples}}

\subsection*{Function}

\code{ScanCtl} holds what a program sets and reads of a scanout from
memory (issue 151): where the next frame starts, whether the screen
shows the scanout or the video peripheral's own framebuffer, whether
lines are asked for one row ahead or two (issue 1523), and whether a
column was ever shown before its word arrived, and whether a line
asked for never came, with that line's address (issue 1197). It also
reads back the longest a line took to come whole, in pixels, and how
many lines had not come whole when their rows began (issue 1523).
Its four outputs are registers, for \code{LinePair} and
\code{VidMux} to read.

What it reads comes back from \code{LinePair} over a channel, as one
\code{ScanState}, rather than on wires: the underflow bit, the stuck
bit, the address, the longest line and the count of late lines.
\code{ScanTap}, which this sheet covers too, puts them on it: it reads
\code{LinePair}'s five wires after \code{LinePair} has driven them and
sends them whenever the channel has room. \code{ScanCtl} drives
\code{clear}, which \code{LinePair} reads, so it runs before
\code{LinePair}, and a wire back would be read a step stale in the run
while the netlist reads it fresh. A channel commits at the end of the
step whichever unit ran first.

\subsection*{Register map}

The unit decodes the word from address bits 4 to 2. In
\code{ScanVideo} it sits above \code{1 << SCAN\_BIT}, \code{0x80}, of
the video slot, so on the board its words are at \code{0x3280} to
\code{0x3298}, the HAL's \code{map::SCAN}.

\dsregs{ScanCtl}

\code{ctrl} reads back as written, both its bits, so a host may set
one and keep the other: the HAL's \code{show} and \code{two\_ahead}
each read \code{ctrl} and write it back with their own bit changed.
A write to \code{clear} of any value clears the underflow and stuck
bits, the longest line and the count of late lines; the word itself is
not kept. \code{stuck\_at} is the first line that did not come since
the stuck bit was last clear, and a clear leaves it until another line
does not come.

\dstables{ScanCtl}

\subsection*{Behaviour}

\begin{itemize}
\item \code{base}, \code{mode}, \code{clear} and \code{two} are driven
from the registers \code{fbase}, \code{show}, \code{clr} and
\code{ahead}; \code{clr} is high for the one cycle after a write to
\code{clear}.
\item \code{mode} also starts the fetch: in \code{ScanVideo} it is
\code{LinePair}'s \code{show}, so no line is read from memory until
the show bit is set (issue 1178).
\item \code{two} is \code{LinePair}'s \code{two}, which the pair takes
at the vertical sync, so the ahead bit counts from the next frame.
\item The underflow bit is \code{seen}, the stuck bit
\code{seen\_stuck}, the address \code{seen\_at}, the longest line
\code{seen\_worst} and the count \code{seen\_lates}, the last values
the channel brought. They can read the old values for a few cycles
after a clear, the values already in the channel, the same in the run
and the netlist.
\item Reads and writes are answered the cycle they are taken.
\end{itemize}

\subsection*{Verification}

\begin{itemize}
\item \code{ex\_scanvideo} sets the base and the scanout bit, reads the
underflow bit clear after a clean frame, set after a frame whose
memory stalls for a line, and clear again after a write to
\code{clear}. \code{//docs:scanvideo\_sim\_scanvideo\_tb\_test} and
\code{//docs:scanvideo\_vsim\_test} simulate it inside
\code{ScanVideo}'s netlist against that run.
\item \code{//lib/parts:parts\_test}'s
\code{ctrl\_reads\_back\_as\_written} writes \code{ctrl} with each
combination of its two bits, reads each back, and sees \code{mode}
and \code{two} follow them.
\item \code{//cpu/vreteno:board\_test}'s
\code{a\_line\_that\_never\_comes\_is\_stuck} points the scanout's
\code{LinePair} at memory that never answers and sees its stuck bit
set with that line's address.
\end{itemize}

\subsection*{Used by}

\code{ScanVideo}. On the board, \code{cpu/vreteno/rust/scanprobe.rs}
uses it through the HAL's \code{Scan}, and says \code{scan stuck} with
the address when the bit is set. \code{cpu/vreteno/rust/ico\_hdmi.rs}
prints the longest line and the late lines each time it reports, then
clears them, and built with \code{ahead} it sets the ahead bit.

\subsection*{Limits}

One base, taken at each vertical sync; there is no interrupt, so a
program reads the underflow and stuck bits.
